% Generated by publication/build_sections.py; do not edit by hand.
\paragraph{Informal verdict.}
\textbf{A --- Complete and apparently correct.} Confidence: \textbf{high (0.96)}.

\paragraph{Lean coverage.}
Generation 0 faithfully defines the half-open A\_n, measurable ENNReal step function f, and actual I/J lintegrals on Real (equivalent to [0,infinity) because f vanishes below 1). It unconditionally proves rational divergence, sqrt\(2\) irrationality/J-finiteness, and Liouville positivity/irrationality/J-divergence. It contains no exact general overlap or double-series/convergence-criterion layer.

\subsection{Audited informal answer}
Complete solution
\begin{enumerate}
\item Faithful restatement
\end{enumerate}
For every positive integer \(n\), let

\[
A_n=[n,n+n^{-3}),
\]

and define \(f:\mathbb R_+\to\mathbb R_+\) by

\[
f(t)=
\begin{cases}
n,&t\in A_n,\\
0,&t\notin \displaystyle\bigcup_{n\ge 1}A_n.
\end{cases}
\]

For \(a>0\), set

\[
J(a)=\int_0^\infty t\,f(t)f(t/a)\,dt,
\qquad
D=\{a>0:J(a)<\infty\}.
\]

The problem asks:


prove that, for every positive rational \(a\),

\[
\int_0^\infty f(t)f(t/a)\,dt=\infty,
\]

and hence \(J(a)=\infty\);


exhibit irrational \(a\) for which \(J(a)<\infty\);


determine whether there exists an irrational \(a\) for which \(J(a)=\infty\). Luc Pommeret+1


There is no apparent typo. The notation \([x,y[\) in the PDF means \([x,y)\). Whether \(\mathbb R_+\) includes \(0\) is immaterial here, since \(f\) vanishes below \(1\).
We shall prove substantially more: an exact convergence criterion in terms of one-sided rational approximations of order \(n^{-4}\).

\begin{enumerate}
\item Preliminary facts
\end{enumerate}
The intervals \(A_n\) are pairwise disjoint, because

\[
n+n^{-3}\le n+1,
\]

with equality only for \(n=1\), where the right endpoint is excluded.
For later use, define also

\[
I(a):=\int_0^\infty f(t)f(t/a)\,dt.
\]

All integrals below are understood as extended nonnegative Lebesgue integrals. Thus Tonelli's theorem is always applicable.
Verification of the averaged integrability stated in the problem
First,

\[
\int_0^\infty f(t)\,dt
=\sum_{n=1}^\infty n\,|A_n|
=\sum_{n=1}^\infty \frac1{n^2}<\infty.
\]

For \(R>0\), Tonelli's theorem and the substitution \(s=t/u\) give

\[
\begin{aligned}
\int_0^R u^{-2}J(u)\,du
&=
\int_0^\infty t f(t)
   \left(\int_0^R u^{-2}f(t/u)\,du\right)dt\\
&=
\int_0^\infty t f(t)
   \left(\frac1t\int_{t/R}^\infty f(s)\,ds\right)dt\\
&=
\int_0^\infty f(t)\int_{t/R}^\infty f(s)\,ds\,dt\\
&\le
\left(\int_0^\infty f\right)^2<\infty.
\end{aligned}
\]

This justifies the contextual assertion that \(J(a)<\infty\) for almost every \(a>0\).

\begin{enumerate}
\item Exact overlap formula
\end{enumerate}
Fix \(a>0\). For \(m,n\ge1\), put

\[
E_{m,n}(a):=A_m\cap aA_n,
\]

where

\[
aA_n=[an,an+a n^{-3}).
\]

On \(E_{m,n}(a)\),

\[
f(t)=m,\qquad f(t/a)=n.
\]

Since both families \((A_m)\) and \((aA_n)\) are pairwise disjoint, the sets \(E_{m,n}(a)\) are pairwise disjoint. Therefore

\[
I(a)=\sum_{m,n\ge1}mn\,|E_{m,n}(a)|
\tag{3.1}
\]

and

\[
J(a)=\sum_{m,n\ge1}mn\int_{E_{m,n}(a)}t\,dt.
\tag{3.2}
\]

Let

\[
r=\frac mn,
\qquad
x_+=\max(x,0).
\]

Define

\[
\Delta_{m,n}(a)=
\begin{cases}
\displaystyle
\left(
\min\left\{a,\,
r^{-3}-n^4(a-r)\right\}
\right)_+,
& r\le a,\\[3mm]
\displaystyle
\left(
\min\left\{r^{-3},\,
a-n^4(r-a)\right\}
\right)_+,
& r>a.
\end{cases}
\tag{3.3}
\]

Then

\[
|E_{m,n}(a)|=n^{-3}\Delta_{m,n}(a).
\tag{3.4}
\]

Indeed, suppose first that \(r\le a\), so \(m\le an\). The later left endpoint is \(an\), and hence

\[
\begin{aligned}
|E_{m,n}(a)|
&=
\left(
\min\{m+m^{-3},an+a n^{-3}\}-an
\right)_+\\
&=
\left(
\min\{m^{-3}-(an-m),a n^{-3}\}
\right)_+\\
&=
n^{-3}
\left(
\min\{r^{-3}-n^4(a-r),a\}
\right)_+.
\end{aligned}
\]

The case \(r>a\) is symmetric.
When \(E_{m,n}(a)\neq\varnothing\), it is an interval with left endpoint

\[
\max(m,an)=n\max(r,a)
\]

and length \(n^{-3}\Delta_{m,n}(a)\). Consequently,

\[
\boxed{
I(a)=
\sum_{m,n\ge1}
\frac{r}{n}\,\Delta_{m,n}(a)
}
\tag{3.5}
\]

and

\[
\boxed{
J(a)=
\sum_{m,n\ge1}
\left[
r\max(r,a)\Delta_{m,n}(a)
+
\frac{r}{2n^4}\Delta_{m,n}(a)^2
\right].
}
\tag{3.6}
\]

These formulas are exact.
A useful convergence criterion
If \(E_{m,n}(a)\neq\varnothing\), then for some \(t\in E_{m,n}(a)\),

\[
an\le t<m+m^{-3}\le m+1
\]

and

\[
m\le t<an+a n^{-3}\le an+a.
\]

Thus

\[
an-1<m<an+a.
\tag{3.7}
\]

In particular, for every fixed \(n\), only boundedly many \(m\)'s can occur, with a bound depending only on \(a\). Moreover, for all sufficiently large \(n\),

\[
\frac a2<\frac mn<2a
\qquad
\text{whenever }\Delta_{m,n}(a)>0.
\tag{3.8}
\]

Also,

\[
0\le \Delta_{m,n}(a)\le \max\{a,a^{-3}\}.
\tag{3.9}
\]

The second sum in (3.6),

\[
\sum_{m,n}\frac{r}{2n^4}\Delta_{m,n}(a)^2,
\]

therefore always converges: for each \(n\) there are only \(O_a(1)\) relevant \(m\)'s, while the summand is \(O_a(n^{-4})\).
On the other hand, by (3.8),

\[
\frac{a^2}{2}
\le r\max(r,a)\le 4a^2
\]

for all sufficiently large relevant pairs. We have therefore proved the exact criterion

\[
\boxed{
J(a)<\infty
\quad\Longleftrightarrow\quad
\sum_{m,n\ge1}\Delta_{m,n}(a)<\infty.
}
\tag{3.10}
\]

The quantity \(\Delta_{m,n}(a)\) measures the residual overlap after a rational \(m/n\) approximates \(a\) on the natural scale \(n^{-4}\).

\begin{enumerate}
\item Part \(a\): every positive rational gives divergence
\end{enumerate}
Let

\[
a=\frac pq,
\qquad p,q\in\mathbb N^*.
\]

For each \(k\ge1\), choose

\[
m=pk,\qquad n=qk.
\]

Then \(m/n=a\), and hence

\[
\Delta_{pk,qk}(a)=\min\{a,a^{-3}\}=:\delta_a>0.
\]

By the exact formula (3.5),

\[
I(a)
\ge
\sum_{k=1}^\infty
\frac{a}{qk}\delta_a
=
\frac{a\delta_a}{q}
\sum_{k=1}^\infty\frac1k
=\infty.
\]

Thus

\[
\boxed{
\int_0^\infty f(t)f(t/a)\,dt=\infty
\quad\text{for every }a\in\mathbb Q_{>0}.
}
\]

Furthermore, wherever \(f(t)f(t/a)>0\), one has \(t\ge1\). Therefore

\[
J(a)
=
\int_0^\infty t f(t)f(t/a)\,dt
\ge
\int_0^\infty f(t)f(t/a)\,dt
=\infty.
\]

This proves part \(a\).
As a direct check, when \(a=1\), every pair \(m=n\) contributes

\[
n^2|A_n|=\frac1n
\]

to \(I(1)\), giving precisely the harmonic-series mechanism.

\begin{enumerate}
\item Part \(b\): irrational examples with finite \(J(a)\)
\end{enumerate}
We first give a general sufficient condition.
Proposition 5.1
Suppose that there exist constants \(c>0\), \(\sigma<4\), and \(N_0\) such that

\[
\left\lvert{}a-\frac mn\right\rvert{}
\ge \frac{c}{n^\sigma}
\tag{5.1}
\]

for every \(n\ge N_0\) and every \(m\ge1\). Then \(J(a)<\infty\).
Proof
Suppose \(\Delta_{m,n}(a)>0\), with \(n\) sufficiently large.
If \(r=m/n\le a\), then by (3.3),

\[
r^{-3}-n^4(a-r)>0.
\]

By (3.8), \(r\ge a/2\), and therefore

\[
0<a-r<r^{-3}n^{-4}
\le 8a^{-3}n^{-4}.
\]

If \(r>a\), positivity of \(\Delta_{m,n}(a)\) gives

\[
0<r-a<a n^{-4}.
\]

Thus in all cases,

\[
\left\lvert{}a-\frac mn\right\rvert{}
<C_a n^{-4},
\qquad
C_a:=\max\{a,8a^{-3}\}.
\tag{5.2}
\]

But (5.1) and (5.2) are incompatible for sufficiently large \(n\), since \(\sigma<4\). Hence only finitely many pairs \((m,n)\) have nonempty overlap. Formula (3.2) then shows that \(J(a)<\infty\). \(\square\)
Explicit examples: quadratic surds
Let \(d\ge2\) be an integer that is not a square, and set

\[
a=\sqrt d.
\]

For \(r=m/n\le2a\),

\[
\begin{aligned}
|a-r|
&=
\frac{|a^2-r^2|}{a+r}\\
&=
\frac{|dn^2-m^2|}{n^2(a+r)}.
\end{aligned}
\]

Since \(d\) is not a square,

\[
dn^2-m^2\neq0,
\]

so its absolute value is at least \(1\). Since \(a+r\le3a\),

\[
\left\lvert{}\sqrt d-\frac mn\right\rvert{}
\ge
\frac1{3\sqrt d\,n^2}.
\tag{5.3}
\]

If \(r>2a\), then \(|a-r|>a\), which is also bounded below by a positive multiple of \(n^{-2}\). Consequently there exists \(c_d>0\) such that

\[
\left\lvert{}\sqrt d-\frac mn\right\rvert{}
\ge \frac{c_d}{n^2}
\]

for all positive integers \(m,n\).
Proposition 5.1 applies with \(\sigma=2\). Therefore

\[
\boxed{
J(\sqrt d)<\infty
\quad
\text{for every nonsquare integer }d\ge2.
}
\]

In particular,

\[
\boxed{J(\sqrt2)<\infty.}
\]

This answers part \(b\), with infinitely many explicit examples.

\begin{enumerate}
\item Part \(c\): an irrational with \(J(a)=\infty\)
\end{enumerate}
We first isolate the general divergence mechanism.
Proposition 6.1
Let \(a>0\). Suppose that there are positive integers \(m_k,n_k\), with \(n_k\to\infty\), such that

\[
n_k^4\left\lvert{}a-\frac{m_k}{n_k}\right\rvert{}\longrightarrow0.
\tag{6.1}
\]

Then

\[
J(a)=\infty.
\]

Proof
Put \(r_k=m_k/n_k\). Condition (6.1) implies \(r_k\to a\).
If \(r_k\le a\), then by (3.3),

\[
\Delta_{m_k,n_k}(a)
=
\left(
\min\left\{
a,\,
r_k^{-3}-n_k^4(a-r_k)
\right\}
\right)_+.
\]

If \(r_k>a\), then

\[
\Delta_{m_k,n_k}(a)
=
\left(
\min\left\{
r_k^{-3},\,
a-n_k^4(r_k-a)
\right\}
\right)_+.
\]

In either case,

\[
\Delta_{m_k,n_k}(a)
\longrightarrow
\min\{a,a^{-3}\}>0.
\]

Thus the series \(\sum_{m,n}\Delta_{m,n}(a)\) contains infinitely many terms bounded below by a fixed positive constant. It diverges, and (3.10) gives \(J(a)=\infty\). \(\square\)
An explicit irrational example
Define

\[
L:=\sum_{j=1}^\infty 10^{-j!}.
\tag{6.2}
\]

For \(k\ge1\), set

\[
q_k=10^{k!},
\qquad
p_k=q_k\sum_{j=1}^k10^{-j!}.
\]

Because \(j!\le k!\), \(p_k\) is a positive integer, and

\[
\frac{p_k}{q_k}=\sum_{j=1}^k10^{-j!}<L.
\]

The tail satisfies

\[
0<L-\frac{p_k}{q_k}
=\sum_{j=k+1}^\infty10^{-j!}
<2\cdot10^{-(k+1)!}.
\tag{6.3}
\]

Indeed, from the first omitted term onward, successive terms decrease by factors bounded above by \(10^{-2}\).
Since

\[
10^{-(k+1)!}
=
q_k^{-(k+1)},
\]

we obtain

\[
q_k^4\left(L-\frac{p_k}{q_k}\right)
<
2q_k^{\,3-k}
=
2\cdot10^{(3-k)k!}
\longrightarrow0.
\tag{6.4}
\]

It remains to verify that \(L\) is irrational. Suppose instead that \(L=r/s\) with positive integers \(r,s\). Since \(p_k/q_k<L\), the two rationals are distinct, and hence

\[
L-\frac{p_k}{q_k}
=
\left\lvert{}\frac rs-\frac{p_k}{q_k}\right\rvert{}
\ge \frac1{s q_k}.
\tag{6.5}
\]

But (6.3) gives

\[
L-\frac{p_k}{q_k}
<2q_k^{-(k+1)}
<\frac1{s q_k}
\]

for all sufficiently large \(k\), a contradiction. Thus \(L\) is irrational.
Applying Proposition 6.1 with \(m_k=p_k\), \(n_k=q_k\), we conclude that

\[
\boxed{
J(L)=\infty,
\qquad
L=\sum_{j=1}^\infty10^{-j!}\notin\mathbb Q.
}
\]

Hence the answer to part \(c\) is yes.

\begin{enumerate}
\item Final theorem
\end{enumerate}
For the function \(f\) of Q885:

\[
\boxed{
\begin{aligned}
&\text{(a) Every }a\in\mathbb Q_{>0}\text{ satisfies }
\int_0^\infty f(t)f(t/a)\,dt=\infty,
\text{ hence }J(a)=\infty.\\[1mm]
&\text{(b) Every }\sqrt d,\ d\ge2\text{ nonsquare, satisfies }J(\sqrt d)<\infty.\\[1mm]
&\text{(c) The explicit irrational }
L=\sum_{j=1}^\infty10^{-j!}
\text{ satisfies }J(L)=\infty.
\end{aligned}
}
\]

More generally, the governing scale is quartic rational approximation:

\[
\left\lvert{}a-\frac mn\right\rvert{}\asymp n^{-4}.
\]

A Diophantine lower bound of order \(n^{-\sigma}\) with \(\sigma<4\) forces only finitely many overlaps and hence finite \(J(a)\), whereas infinitely many approximations satisfying

\[
n^4\left\lvert{}a-\frac mn\right\rvert{}\longrightarrow0
\]

force \(J(a)=\infty\).
Finally, the substitution \(t=as\) yields the useful symmetry

\[
J(a)=a^2J(1/a),
\]

so \(D\) is invariant under \(a\mapsto1/a\). Thus \(1/\sqrt d\) gives further finite examples, while \(1/L\) gives another irrational divergent example.
